%% IEEE Conference Paper — two-column
%% Compile with pdflatex (needs IEEEtran.cls, which ships with most TeX distros).
%%   pdflatex research_paper.tex  (run twice for references/labels)
%% Overleaf: create a project, set "IEEEtran" document class (built in).

\documentclass[conference]{IEEEtran}
\IEEEoverridecommandlockouts

\usepackage{cite}
\usepackage{amsmath,amssymb,amsfonts}
\usepackage{algorithmic}
\usepackage{graphicx}
\usepackage{textcomp}
\usepackage{xcolor}
\usepackage{booktabs}
\usepackage{array}
\def\BibTeX{{\rm B\kern-.05em{\sc i\kern-.025em b}\kern-.08em T\kern-.1667em\lower.7ex\hbox{E}\kern-.125emX}}

\begin{document}

\title{Adaptive Traffic Signal Control at an Isolated Intersection using a
Dueling Double Deep Q-Network with PCU-Weighted State Representation for
Heterogeneous Traffic}

\author{\IEEEauthorblockN{Ali Mehdi Mirza}
\IEEEauthorblockA{\textit{Department of [Your Department]} \\
\textit{[Your Institution]}\\
[City, Country] \\
alimehdi1010@gmail.com}
}

\maketitle

\begin{abstract}
Fixed-time and vehicle-actuated traffic signal controllers are the de-facto
standard at urban intersections, yet both degrade under asymmetric and
time-varying demand because their control logic is not learned from traffic
dynamics. This paper presents an adaptive signal controller based on a
Dueling Double Deep Q-Network (D3QN) that selects signal phases from a
discretized, normalized intersection state and is trained entirely in the
microscopic simulator SUMO through the TraCI interface. We formulate signal
control as a discrete phase-selection Markov Decision Process with a shaped,
bounded reward centred on reducing cumulative vehicle waiting time, and delegate
minimum-green and yellow/all-red clearance enforcement to the environment to
guarantee physically valid, safe signal plans. Evaluated against fixed-time and
actuated baselines on three demand regimes (balanced, asymmetric, and a
mid-episode surge), each replicated over multiple random seeds and reported as
mean $\pm$ standard deviation, the D3QN reduces average waiting time by
75.5--82.6\% relative to fixed-time control and by 18.1--38.0\% relative to a
strong actuated baseline, while matching throughput and reducing estimated
CO\textsubscript{2} emissions by nearly an order of magnitude. We further
address the common homogeneous-fleet assumption by introducing a Passenger Car
Unit (PCU) weighting (IRC~106-1990) into both state and reward, evaluated on a
mixed fleet (car, motorcycle, auto-rickshaw, bus) under SUMO's sub-lane model.
An ablation shows PCU weighting lowers waiting time by 23.9\% and halves
run-to-run variance relative to an identical unweighted agent. All code, models,
and experiments are reproducible via a single containerized command.
\end{abstract}

\begin{IEEEkeywords}
Reinforcement learning, deep Q-network, dueling double DQN, adaptive traffic
signal control, SUMO, TraCI, intelligent transportation systems, passenger car
unit, heterogeneous traffic, queue length, waiting time, Markov decision process.
\end{IEEEkeywords}

\section{Introduction}
Urban road intersections are principal bottlenecks in city traffic networks.
Conventional signal control falls into two families: \emph{fixed-time} control,
which cycles through phases on a pre-computed schedule irrespective of live
demand, and \emph{vehicle-actuated} control, which extends a green phase while
detectors sense arriving vehicles. Both are effective under the conditions for
which they were tuned, but neither \emph{learns} from the traffic it governs, so
both lose efficiency when demand is asymmetric across approaches or varies over
time.

Reinforcement learning (RL) allows an agent to learn a control policy directly
from interaction with the traffic environment, optimizing a long-horizon
objective such as cumulative delay. Deep Q-Networks (DQN)~\cite{mnih2015} and
their successors---Double DQN~\cite{vanhasselt2016} and the Dueling
architecture~\cite{wang2016}---have been applied to signal control. However,
three practical gaps recur: (i) \emph{weak baselines}, where studies compare only
against fixed-time control and omit the stronger actuated baseline; (ii) the
\emph{homogeneous-traffic assumption}, where states count raw vehicles, treating
a motorcycle and a bus as equivalent; and (iii) limited \emph{reproducibility}
from single-run results without seed replication.

This paper addresses all three by benchmarking against both fixed-time and
actuated control, introducing a PCU-weighted state and reward evaluated on a
heterogeneous fleet, and releasing a containerized, multi-seed pipeline.

\subsection{Contributions}
\begin{itemize}
\item A Dueling Double DQN controller for an isolated 4-way intersection, with
the environment (not the agent) enforcing minimum-green and mandatory
yellow/all-red clearance for safety.
\item A shaped, bounded reward with per-component logging for transparent
behaviour analysis.
\item A multi-seed evaluation against fixed-time and actuated baselines across
three demand regimes, reporting waiting time, queue length, throughput, and
estimated emissions.
\item A PCU-weighted state/reward formulation for heterogeneous traffic and an
ablation demonstrating its benefit.
\item A one-command, Docker-based reproducible artifact.
\end{itemize}

\section{Related Work and Research Gaps}
RL-based signal control commonly pairs SUMO~\cite{sumo2018} and TraCI with a
DQN-family agent. Foundational techniques include experience replay and target
networks~\cite{mnih2015}, the Double DQN correction for Q-value
overestimation~\cite{vanhasselt2016}, and the Dueling architecture separating
state value from action advantage~\cite{wang2016}. Open-source interfaces such
as SUMO-RL~\cite{sumorl} standardize the environment API.

The gaps this work fulfils are: \emph{baseline rigour} (we include the actuated
controller as a strong baseline); \emph{heterogeneity} (we embed IRC~106-1990
PCU factors~\cite{irc106} into the learned representation and quantify the
effect); and \emph{reproducibility} (mean~$\pm$~std over seeds and a
containerized pipeline).

\section{Methodology}

\subsection{Simulation Environment}
Experiments use Eclipse SUMO~1.27.1 (headless) controlled via TraCI. The network
is a single 4-way intersection generated with \texttt{netconvert}: four incoming
and four outgoing edges, two lanes each, 200~m length, 13.89~m/s free-flow speed.
Left turns are permissive (gap-acceptance), yielding a compact two-phase
(North--South / East--West) control structure.

\subsection{Demand Scenarios}
\emph{Balanced} ($\sim$300 veh/h/approach, symmetric); \emph{Asymmetric} (heavy
N--S, light E--W); \emph{Surge} (mid-episode demand ramp); and
\emph{Heterogeneous} (mixed fleet under SUMO's SL2015 sub-lane model, lateral
resolution 0.8~m).

\subsection{MDP Formulation}
\textbf{State (71-D, normalized).} Per incoming lane: a 3-cell density vector and
a 3-cell mean-speed vector, per-lane queue length and cumulative waiting time,
a 6-D current-phase one-hot, and a minimum-green-elapsed flag:
$8{\times}3 + 8{\times}3 + 8 + 8 + 6 + 1 = 71$.

\textbf{Action ($|A|{=}2$).} $a{=}0$ requests N--S green; $a{=}1$ requests E--W
green; decisions every 5~s.

\textbf{Safety wrapper.} The environment enforces a 10~s minimum green and inserts
3~s yellow $+$ 2~s all-red clearance on every axis change.

\textbf{Reward.}
\begin{equation}
\begin{aligned}
r_t = &\;(W_{t-1} - W_t) - 0.5\,\mathbb{1}[a_t \neq a_{t-1}] - 0.1\!\sum_l q_{l,t},\\
r_t \leftarrow &\; \mathrm{clip}(r_t,\,-10,\,+10),
\end{aligned}
\end{equation}
where $W_t$ is total cumulative waiting time and $q_{l,t}$ the queue on lane $l$.
Components are logged separately.

\subsection{Agent: Dueling Double DQN}
A shared trunk (two hidden layers of 128 ReLU units) branches into value $V(s)$
and advantage $A(s,a)$ streams, aggregated with mean subtraction:
\begin{equation}
Q(s,a) = V(s) + \Big( A(s,a) - \tfrac{1}{|A|}\textstyle\sum_{a'} A(s,a') \Big).
\end{equation}
Double DQN targets decouple selection from evaluation:
\begin{equation}
y_t = r_t + \gamma\, Q_{\theta^-}\!\big(s_{t+1}, \arg\max_a Q_\theta(s_{t+1},a)\big)(1-d_t).
\end{equation}
Stabilizers: replay buffer (50\,000), hard target update every 500 steps,
$\varepsilon$-greedy decay ($1.0\!\rightarrow\!0.05$), Huber loss, gradient
clipping; Adam ($\eta{=}5\times10^{-4}$), $\gamma{=}0.99$, batch~64.

\subsection{PCU-Weighted Representation}
For mixed fleets, each vehicle's contribution to the queue features and the
reward's queue penalty is weighted by its PCU factor~\cite{irc106}:
car~$=1.0$, motorcycle~$=0.5$, auto-rickshaw~$=0.8$, bus~$=3.0$. Raw counts are
retained for fair metric reporting.

\subsection{Metrics}
Average waiting time (s, primary), average queue length, throughput (veh/h), and
estimated CO\textsubscript{2} (mg). Each configuration is run over three seeds;
we report mean~$\pm$~std. Training is monitored via the episode-reward curve.

\section{Experimental Setup}
Evaluation episodes span 1200~simulated seconds; agents train for 80 episodes
(homogeneous) and 150 (heterogeneous), with $\varepsilon$-decay matched to the
budget. All runs execute in a Docker image (Ubuntu~22.04 + SUMO + PyTorch~2.14
CPU). The DQN is evaluated greedily ($\varepsilon{=}0$) on seeds $\{1,2,3\}$.

\section{Results}

\begin{table}[t]
\caption{Average Waiting Time (s), mean $\pm$ std over 3 seeds.}
\label{tab:wait}
\centering
\begin{tabular}{@{}lcccc@{}}
\toprule
Scenario & Fixed-time & Actuated & \textbf{D3QN} & $\Delta$ vs fixed \\
\midrule
Balanced   & 56.85$\pm$3.53 & 19.29$\pm$2.66 & \textbf{11.97$\pm$0.76} & $-$78.9\% \\
Asymmetric & 64.64$\pm$1.11 & 23.82$\pm$0.17 & \textbf{15.82$\pm$0.62} & $-$75.5\% \\
Surge      & 34.88$\pm$0.17 & \ \,7.42$\pm$0.94 & \ \,\textbf{6.08$\pm$1.77} & $-$82.6\% \\
\bottomrule
\end{tabular}
\end{table}

\begin{table}[t]
\caption{Average Queue Length (halting vehicles).}
\label{tab:queue}
\centering
\begin{tabular}{@{}lccc@{}}
\toprule
Scenario & Fixed-time & Actuated & \textbf{D3QN} \\
\midrule
Balanced   & 4.14$\pm$0.13 & 2.29$\pm$0.09 & \textbf{1.93$\pm$0.10} \\
Asymmetric & 4.71$\pm$0.05 & 2.45$\pm$0.06 & \textbf{2.11$\pm$0.07} \\
Surge      & 2.71$\pm$0.00 & 1.16$\pm$0.09 & \textbf{0.96$\pm$0.07} \\
\bottomrule
\end{tabular}
\end{table}

\begin{table}[t]
\caption{Throughput (veh/h) --- matched across controllers.}
\label{tab:thru}
\centering
\begin{tabular}{@{}lccc@{}}
\toprule
Scenario & Fixed-time & Actuated & D3QN \\
\midrule
Balanced   & 1168 & 1170 & 1170 \\
Asymmetric & 1356 & 1360 & 1365 \\
Surge      & \ \,786 & \ \,784 & \ \,784 \\
\bottomrule
\end{tabular}
\end{table}

\begin{table}[t]
\caption{Estimated CO\textsubscript{2} ($\times10^6$ mg per episode).}
\label{tab:co2}
\centering
\begin{tabular}{@{}lccc@{}}
\toprule
Scenario & Fixed-time & Actuated & \textbf{D3QN} \\
\midrule
Balanced   & 19.4 & 15.5 & \textbf{2.1} \\
Asymmetric & 23.3 & 17.8 & \textbf{2.6} \\
Surge      & 12.6 & \ \,9.5 & \textbf{1.5} \\
\bottomrule
\end{tabular}
\end{table}

\subsection{Main Comparison (Homogeneous Demand)}
Table~\ref{tab:wait} shows the D3QN reduces average waiting time by 75.5--82.6\%
versus fixed-time and 18.1--38.0\% versus actuated control. Queue lengths
(Table~\ref{tab:queue}) are correspondingly lowest for the D3QN, while throughput
(Table~\ref{tab:thru}) is matched: the improvement is in service efficiency, not
vehicles cleared. Estimated CO\textsubscript{2} (Table~\ref{tab:co2}) drops by
roughly an order of magnitude; as simulator estimates these are indicative.

\subsection{Training Behaviour}
Across scenarios the episode reward increases and stabilizes once $\varepsilon$
reaches its floor, with bounded loss and no divergence, confirming stable
learning under the Double-DQN/dueling stabilizers and bounded reward.

\subsection{Heterogeneous Traffic: PCU Ablation}
\begin{table}[t]
\caption{Heterogeneous Traffic --- Avg.\ Waiting Time (s), 3 seeds.}
\label{tab:pcu}
\centering
\begin{tabular}{@{}lcc@{}}
\toprule
Controller & Avg wait (s) & $\Delta$ vs fixed-time \\
\midrule
Fixed-time & 56.11$\pm$0.95 & --- \\
Actuated & 18.53$\pm$1.49 & $-$67.0\% \\
D3QN (no PCU) & 29.38$\pm$4.85 & $-$47.6\% \\
\textbf{D3QN (PCU-weighted)} & \textbf{22.36$\pm$1.95} & \textbf{$-$60.1\%} \\
\bottomrule
\end{tabular}
\end{table}
Table~\ref{tab:pcu} shows PCU weighting improves the D3QN by 23.9\% in waiting
time over the identical unweighted agent (22.36 vs 29.38~s) and reduces variance
from $\pm4.85$ to $\pm1.95$. At the training budget used, the PCU-weighted agent
does not yet surpass the actuated baseline on heterogeneous traffic; closing this
gap is future work.

\subsection{On Scientific Process}
An initial short heterogeneous run (60 episodes, PCU in state only) produced the
opposite result. Diagnosis attributed this to under-training and an unchanged
reward objective. A corrected protocol (PCU in both state and reward, 150
episodes, retuned $\varepsilon$-decay) reversed the finding to the positive
result in Table~\ref{tab:pcu}. Both are reported for transparency.

\section{Discussion}
A compact D3QN learns a signal policy that outperforms fixed-time and a
well-configured actuated controller on delay and queueing without sacrificing
throughput. The safety wrapper guarantees valid minimum-green and clearance
intervals. The PCU ablation supports the argument that vehicle-class-aware state
and reward yield a better, more stable policy for heterogeneous traffic than
uniform counting.

\section{Limitations}
Single isolated intersection; simulation-only with permissive left turns and a
discretized state; limited training budget (1200~s episodes, 80--150 episodes);
emissions are simulator estimates; and the learned agent does not yet beat the
actuated baseline on heterogeneous traffic at the current budget.

\section{Future Work}
Protected left-turn phases and larger action sets; multi-intersection
coordination via multi-agent RL; a single generalist policy across regimes;
richer state encoders and longer training for heterogeneous traffic; prioritized
replay and distributional value estimation.

\section{Conclusion}
We presented a Dueling Double DQN signal controller trained in SUMO and evaluated
against fixed-time and actuated baselines across three demand regimes with
multi-seed statistics. It reduces average waiting time by 75.5--82.6\% versus
fixed-time and up to 38\% versus actuated control, with matched throughput and
substantially lower estimated emissions. A PCU-weighted extension improves the
agent by 23.9\% over an unweighted counterpart on heterogeneous traffic. The
full pipeline is containerized and reproducible with a single command.

\section*{Reproducibility}
Code, models, and scripts:
\texttt{github.com/ali0786mehdi/Intelligent-Traffic-Signal-Controller}.
Run \texttt{make all} (Linux/Mac) or \texttt{.\textbackslash run.ps1 all}
(Windows) inside Docker to reproduce all tables and figures.

\begin{thebibliography}{00}
\bibitem{mnih2015} V. Mnih \emph{et al.}, ``Human-level control through deep
reinforcement learning,'' \emph{Nature}, vol.~518, pp.~529--533, 2015.
\bibitem{vanhasselt2016} H. van Hasselt, A. Guez, and D. Silver, ``Deep
reinforcement learning with double Q-learning,'' in \emph{Proc. AAAI}, 2016.
\bibitem{wang2016} Z. Wang \emph{et al.}, ``Dueling network architectures for
deep reinforcement learning,'' in \emph{Proc. ICML}, 2016.
\bibitem{sumo2018} P. A. Lopez \emph{et al.}, ``Microscopic traffic simulation
using SUMO,'' in \emph{Proc. IEEE ITSC}, 2018.
\bibitem{irc106} Indian Roads Congress, \emph{IRC:106-1990 --- Guidelines for
Capacity of Urban Roads in Plain Areas}, 1990.
\bibitem{sumorl} L. N. Alegre, ``SUMO-RL,'' open-source Gymnasium interface for
SUMO, 2019.
\end{thebibliography}

\end{document}
